\documentclass[11pt]{article}

% --- Paquets de base ---
\usepackage[utf8]{inputenc}
\usepackage[T1]{fontenc}
\usepackage[french]{babel}
\usepackage{amsmath, amsfonts, amssymb}
\usepackage{geometry}
\geometry{a4paper, margin=1in}

% --- Mise en page et Style ---
\usepackage{titlesec}
\usepackage{enumitem}
\usepackage{booktabs}
\usepackage{longtable}
\usepackage{xcolor}
\usepackage{tcolorbox}
\usepackage{tikz}
\usepackage{hyperref}
\usepackage{caption}
\usepackage{graphicx}
\usetikzlibrary{shapes.geometric, arrows.meta, positioning, calc}

% --- Macros Mathématiques ---
\newcommand{\E}{\mathbb{E}}
\newcommand{\Prob}{\mathbb{P}}
\newcommand{\R}{\mathbb{R}}
\newcommand{\N}{\mathbb{N}}
\newcommand{\Var}{\mathrm{Var}}
\newcommand{\Cov}{\mathrm{Cov}}
\newcommand{\calN}{\mathcal{N}}
\newcommand{\calF}{\mathcal{F}}
\newcommand{\ind}{\mathbf{1}}

% --- Boîtes de définition ---
\tcbuselibrary{skins, breakable}
\newtcolorbox{definitionbox}[1][]{
  colback=gray!8, colframe=black!60,
  fonttitle=\bfseries, title=#1,
  sharp corners, boxrule=0.6pt, breakable
}
\newtcolorbox{noveltybox}[1][]{
  colback=blue!4, colframe=blue!50,
  fonttitle=\bfseries, title=#1,
  sharp corners, boxrule=0.8pt, breakable
}
\newtcolorbox{resourcebox}[1][]{
  colback=green!4, colframe=green!50,
  fonttitle=\bfseries, title=#1,
  sharp corners, boxrule=0.6pt, breakable
}
\newtcolorbox{resultbox}[1][]{
  colback=orange!6, colframe=orange!60!black,
  fonttitle=\bfseries, title=#1,
  sharp corners, boxrule=0.8pt, breakable
}
\newtcolorbox{diagnosticbox}[1][]{
  colback=red!4, colframe=red!50,
  fonttitle=\bfseries, title=#1,
  sharp corners, boxrule=0.8pt, breakable
}

% --- Configuration des Titres ---
\titleformat{\section}
  {\Large\bfseries\color{black}}{\thesection}{1em}{}[\titlerule]
\titleformat{\subsection}
  {\large\bfseries\color{black}}{\thesubsection}{1em}{}
\titleformat{\subsubsection}
  {\bfseries\color{black}}{\thesubsubsection}{1em}{}

% --- Paramètres de paragraphe ---
\setlength{\parskip}{0.5em}
\setlength{\parindent}{0pt}

\title{Interprétabilité automatique de mails clients par Sparse Autoencoders\\
\large Application à la correspondance client d'EDF}
\author{Grégoire Pelletier}
\date{}

% =============================================================================
\begin{document}

% ─────────────────────────────────────────────────────────────────────────────
% PAGE DE GARDE
% ─────────────────────────────────────────────────────────────────────────────
\begin{titlepage}
\centering
\vspace*{1cm}
{\large Université Paris-Saclay --- Faculté des Sciences d'Orsay\par}
{\large Master 2 Mathématiques et Intelligence Artificielle\par}
\vspace{2cm}
{\Huge\bfseries Rapport de stage\par}
\vspace{0.5cm}
{\LARGE Interprétabilité automatique de mails clients par\\ Sparse Autoencoders\par}
\vspace{0.3cm}
{\Large\itshape Application à l'analyse interprétable de la\\ correspondance client d'EDF\par}
\vspace{2.5cm}

\begin{tabular}{ll}
\textbf{Auteur} & Grégoire Pelletier \\[0.3em]
\textbf{Entreprise d'accueil} & EDF R\&D --- Département SEQUOIA, groupe SOAD \\[0.3em]
\textbf{Lieu} & EDF Lab Paris-Saclay, 7 boulevard Gaspard Monge, 91120 Palaiseau \\[0.3em]
\textbf{Maître de stage} & Philippe Suignard (département SEQUOIA) \\[0.3em]
\textbf{Durée du stage} & 6 mois --- mars 2026 à septembre 2026 \\[0.3em]
\end{tabular}

\vfill
{\large \today\par}
\end{titlepage}

\newpage
\section*{Remerciements}
\addcontentsline{toc}{section}{Remerciements}

Je tiens à remercier Philippe Suignard, Claire Del Gatto et l'ensemble de l'équipe du groupe
Statistiques et Outils d'Aide à la Décision (SOAD) du département SEQUOIA
d'EDF R\&D pour leur accueil, leur disponibilité et les échanges réguliers qui
ont accompagné ce stage.

\newpage
\section*{Résumé}
\addcontentsline{toc}{section}{Résumé}

Ce stage porte sur l'explicabilité automatique de mails clients d'EDF à
l'aide de Sparse Autoencoders (SAE), combinant un SAE préentraîné à grande
échelle (GemmaScope-2, sur les activations de Gemma-3-12B-it) étendu par un
second SAE entraîné spécifiquement sur le domaine (architecture à cœur gelé,
structurellement équivalente à SAE Boost), et un second pipeline indépendant
fondé sur des embeddings de phrase (F2LLM-v2, bge-m3). Le pipeline initial,
fonctionnel de bout en bout, présentait un taux de succès faible (20\%) au
test d'auto-interprétation des features propres au domaine. Une démarche de
diagnostic par ablation contrôlée a établi que ce taux n'était pas limité par
le volume d'entraînement, mais par un défaut de conception du corpus
d'entraînement (uniquement générique, sans mails du domaine cible) ---
corrigé, il porte le taux d'interprétabilité à 45,3\%.

Une campagne d'ablations exhaustive (plus de 20 configurations testées : largeur
du SAE, capacité et parcimonie de l'extension, volume de tokens, graine
d'entraînement, taille du modèle extracteur/juge, dimension d'embedding,
backbone de Pipeline 2) montre qu'\textbf{aucun hyperparamètre du SAE ne modifie
significativement le taux d'interprétabilité} une fois le domaine du corpus
corrigé --- à l'exception d'un unique levier : \textbf{l'échelle du modèle
extracteur/juge}, qui produit un effet dose-réponse net et hautement
significatif (12,0\% à 1B paramètres, 28,0\% à 4B, 45,3\% à 12B ; test de
tendance de Cochran-Armitage, $p \approx 1{,}6\times10^{-10}$). Des tests
complémentaires (fidélité et plausibilité de l'explication document-level,
sanity check contre un décodeur figé aléatoire, robustesse du protocole de
jugement, biais multilingue, fidélité du steering, évaluation quantitative du
retrieval) complètent la validation du système, avec un audit rétroactif de
la méthodologie statistique employée (tests appariés, correction multi-tests).

\textbf{Mots-clés} : Sparse Autoencoders, interprétabilité mécaniste,
GemmaScope, grands modèles de langage, explicabilité, traitement automatique
des mails clients, auto-interprétation par juge LLM.

\section*{Abstract}
\addcontentsline{toc}{section}{Abstract}

This internship addresses automatic explainability of customer emails at EDF
using Sparse Autoencoders (SAE), combining a large pretrained SAE
(GemmaScope-2, on Gemma-3-12B-it activations) extended by a second SAE
trained specifically on the target domain (frozen-core architecture,
structurally equivalent to SAE Boost), alongside an independent
sentence-embedding-based pipeline (F2LLM-v2, bge-m3). The initial end-to-end
pipeline showed a low success rate (20\%) on the domain-specific feature
auto-interpretation test. A controlled-ablation diagnostic established that
this was not a training-volume limitation but a training-corpus design flaw
(generic text only, no target-domain emails) --- once fixed, the
interpretability rate rose to 45.3\%.

An exhaustive ablation campaign (20+ configurations: SAE width, extension
capacity and sparsity, token volume, training seed, extractor/judge model
scale, embedding dimension, Pipeline-2 backbone) shows that \textbf{no SAE
hyperparameter significantly changes the interpretability rate} once the
corpus domain is fixed --- except for a single lever: \textbf{the scale of the
extractor/judge model}, which produces a clean, highly significant
dose-response effect (12.0\% at 1B parameters, 28.0\% at 4B, 45.3\% at 12B;
Cochran-Armitage trend test, $p \approx 1.6\times10^{-10}$). Complementary
tests (document-level explanation fidelity and plausibility, a random-decoder
sanity check, judge robustness, multilingual bias, steering fidelity,
quantitative retrieval evaluation) complete the system's validation, together
with a retroactive audit of the statistical methodology used (paired tests,
multiple-comparison correction).

\textbf{Keywords}: Sparse Autoencoders, mechanistic interpretability,
GemmaScope, large language models, explainability, customer email analysis,
LLM auto-interpretation.

\newpage
\tableofcontents
\newpage
\listoffigures
\listoftables
\newpage

\section*{Glossaire}
\addcontentsline{toc}{section}{Glossaire}
\begin{description}[leftmargin=3.2cm, style=nextline]
  \item[SAE] Sparse Autoencoder (autoencodeur creux) : réseau qui reconstruit un
    vecteur d'entrée à partir d'une représentation intermédiaire parcimonieuse.
  \item[LLM] Large Language Model (grand modèle de langage).
  \item[GemmaScope] Famille de SAE préentraînés par Google DeepMind sur les
    activations des modèles Gemma/Gemma-2/Gemma-3 \cite{lieberum2024gemmascope}.
  \item[L0] Pseudo-norme comptant le nombre de coordonnées non nulles d'un
    vecteur --- mesure la parcimonie effective d'un code SAE.
  \item[NMSE] Normalized Mean Squared Error --- erreur de reconstruction
    normalisée par la variance du signal d'origine.
  \item[FVE] Fraction de Variance Expliquée (parfois notée $R^2$).
  \item[MRL] Matryoshka Representation Learning --- apprentissage d'embeddings
    dont les préfixes de dimensions sont eux-mêmes des embeddings valides.
  \item[BM25] Fonction de pondération classique en recherche d'information,
    ici appliquée au vocabulaire latent d'un SAE (Latent Terms).
  \item[odd-one-out] Protocole d'auto-interprétation : un juge LLM doit
    identifier l'intrus parmi des extraits activant fortement une même feature.
  \item[McNemar] Test statistique apparié pour comparer deux proportions
    mesurées sur les mêmes unités d'observation sous deux conditions.
  \item[Cochran-Armitage] Test de tendance pour une variable binaire face à un
    facteur ordonné à plusieurs niveaux.
\end{description}

\newpage
% ─────────────────────────────────────────────────────────────────────────────
\section{Introduction et contexte}
% ─────────────────────────────────────────────────────────────────────────────

\subsection{Contexte industriel}

Ce stage s'est déroulé au sein d'EDF R\&D, dans le département SEQUOIA
(Services, Économie, Outils Innovants et IA), rattaché au groupe SOAD
(Statistiques et Outils d'Aide à la Décision), sur le site d'EDF Lab
Paris-Saclay. La mission initiale, formulée par l'offre de stage
(\emph{cf.}~annexe des sources), porte sur l'explicabilité des
\textbf{embeddings de texte} : des vecteurs denses obtenus par transformation
d'un document textuel, aujourd'hui massivement utilisés (classification,
recherche de documents similaires) mais dont la production reste, par
construction, une boîte noire. Les Sparse Autoencoders (SAE), apparus
récemment en interprétabilité mécaniste, sont identifiés par l'offre de stage
comme une piste prometteuse pour lever ce verrou.

Le cas d'usage retenu pour ancrer cette investigation méthodologique est
l'analyse de mails clients reçus par EDF : un volume important de
correspondance dont le traitement (priorisation, orientation vers le bon
service, détection des réclamations et des urgences) est aujourd'hui coûteux à
automatiser de façon à la fois efficace et interprétable --- un préalable
important dans un contexte réglementé où une décision automatisée doit
pouvoir être justifiée.

\subsection{Fondements théoriques : polysémanticité et superposition}

Les grands modèles de langage offrent une capacité de compréhension fine du
texte, mais leurs représentations internes restent largement opaques : une
même dimension de leur espace d'activation encode typiquement plusieurs
concepts sémantiques non corrélés, un phénomène appelé
\textbf{polysémanticité}. Il découle de l'hypothèse de \textbf{superposition}
\cite{elhage2022superposition} : contraints par une dimensionnalité fixe, les
réseaux de neurones compressent davantage de caractéristiques (\emph{features})
que de dimensions disponibles, en exploitant des directions quasi-orthogonales
de leur espace latent.

Cette question --- dans quelle mesure un LLM « comprend » réellement le texte
qu'il traite plutôt que d'en imiter la surface statistique --- est elle-même
débattue ; Beckmann \& Queloz \cite{beckmannqueloz2026} soutiennent que les avancées récentes de
l'interprétabilité mécanique rendent la position purement sceptique de moins
en moins tenable, à condition d'articuler ces résultats à un cadre théorique de
la compréhension. Les Sparse Autoencoders proposent précisément de réapprendre
une représentation parcimonieuse et de plus haute dimension dans laquelle
chaque direction correspond, dans l'idéal, à un concept isolé et nommable en
langage naturel \cite{bricken2023monosemanticity, cunningham2023sparse}.

\subsection{Objectifs du stage}

L'énoncé initial fixe l'ambition de construire une plateforme d'analyse de
mails permettant : l'indexation et la recherche par concept, le clustering
interprétable de documents, la détection d'urgence et d'intention, la
comparaison de corpus (diffing), la visualisation des concepts activés, le
retrieval par propriétés, et l'explication des décisions prises en aval de ces
représentations --- avec la contrainte explicite de réutiliser au maximum
l'existant (SAELens, GemmaScope, la littérature sur l'interprétation
d'embeddings documentaires) plutôt que de réimplémenter, et de documenter
systématiquement les écarts avec ces références lorsqu'une réimplémentation
partielle s'avère nécessaire.

Une phase préliminaire de recherche bibliographique (état de l'art détaillé au
chapitre~\ref{sec:etat-art}) a permis de dresser un panorama des architectures
SAE et des modèles d'embeddings candidats, et de formuler plusieurs pistes de
contribution méthodologique (architecture hybride à cœur gelé, pipeline
micro/macro, dimensionnement des données). La suite du stage a consisté à
implémenter, tester empiriquement et, le cas échéant, corriger ces
hypothèses de travail sur des données réelles du domaine.

\begin{definitionbox}[Précision importante sur la nature du corpus]
Le corpus d'emails utilisé (\texttt{Mails.tsv}), désigné dans ce rapport par
commodité comme les mails \og{}originaux\fg{}, n'est \textbf{pas} de la
correspondance client authentique : il s'agit d'un jeu de données déjà
synthétique, produit par un travail antérieur du laboratoire EDF R\&D (non
réalisé pendant ce stage). Les variantes \emph{augmentées} sont, elles,
générées pendant ce stage (Gemma-3-12B-it) à partir de ces mails
\og{}originaux\fg{}. Aucune des deux couches du corpus n'est donc constituée
de données réelles au sens strict --- le terme \og{}original\fg{} désigne
uniquement leur statut de donnée d'entrée par rapport aux variantes qui en
dérivent, pas leur authenticité.
\end{definitionbox}

\subsection{Démarche et plan du rapport}

Le stage a suivi une progression en cinq grandes phases :

\begin{enumerate}
  \item \textbf{Audit et fiabilisation} du pipeline existant (chargement des
    SAE préentraînés, récupération des labels Neuronpedia, précision
    numérique, robustesse) et validation de bout en bout sur un modèle réduit
    avant tout passage à l'échelle.
  \item \textbf{Diagnostic et correction} du problème central identifié en
    début de stage : un taux de succès très faible (20\%) du protocole
    d'auto-interprétation, déterminé par ablation contrôlée comme un défaut de
    conception du corpus, non un manque de volume d'entraînement.
  \item \textbf{Relecture critique face à la littérature de référence} pour
    identifier les écarts méthodologiques restants (retrieval, clustering,
    protocole de labellisation) et les corriger ou les documenter comme
    limites assumées.
  \item \textbf{Mise à l'échelle et consolidation} : tests de qualité de
    l'explication document-level, protocole d'évaluation sous conditions
    fixées, dashboard interactif, première ablation de mise à l'échelle
    (largeur, époques, nombre de features labellisées).
  \item \textbf{Campagne d'ablations exhaustive} couvrant systématiquement les
    modèles et paramètres restants (échelle du modèle extracteur/juge,
    capacité et parcimonie de l'extension isolées, dimension d'embedding,
    backbones alternatifs), avec un audit de la rigueur statistique des
    comparaisons effectuées.
\end{enumerate}

Le chapitre~\ref{sec:etat-art} positionne le projet par rapport à l'état de
l'art. Le chapitre~\ref{sec:contribution} décrit l'architecture et la
méthodologie mises en œuvre. Le chapitre~\ref{sec:resultats} présente
l'ensemble des résultats expérimentaux, cœur scientifique du rapport. Le
chapitre~\ref{sec:discussion} discute les limites, les perspectives et conclut.

\newpage
% ─────────────────────────────────────────────────────────────────────────────
\section{Méthodes, algorithmes et état de l'art}
\label{sec:etat-art}
% ─────────────────────────────────────────────────────────────────────────────

\subsection{Fondements théoriques de l'apprentissage de dictionnaires}

L'hypothèse de la représentation linéaire postule que les modèles de langage
représentent les concepts de haut niveau comme des directions linéaires dans
leur espace d'activation \cite{elhage2022superposition}. Soit une distribution
de données $D_x$ sur un espace d'entrée $\mathcal{X}$ (textes), et une
fonction d'activation interne $h : \mathcal{X} \rightarrow \mathbb{R}^{d_\text{model}}$
extraite d'un modèle. L'objectif de l'apprentissage de dictionnaire
parcimonieux est de trouver une matrice de décodage
$W_\text{dec} \in \mathbb{R}^{d_\text{model} \times d_\text{hidden}}$
($d_\text{hidden} \gg d_\text{model}$) et une fonction d'encodage
$f : \mathbb{R}^{d_\text{model}} \rightarrow \mathbb{R}^{d_\text{hidden}}$
produisant un vecteur d'activation latent $z = f(h(x))$ tel que la
reconstruction $\hat{h}(x) = W_\text{dec}\, z + b_\text{dec} \approx h(x)$ soit
fidèle, tandis que $z$ reste très parcimonieux ($\|z\|_0 \ll d_\text{hidden}$).

\begin{definitionbox}[Convention $W_\text{dec}$ --- nota bene]
Dans ce rapport, $W_\text{dec} \in \mathbb{R}^{d_\text{model} \times d_\text{hidden}}$
et chaque \textbf{colonne} $W_\text{dec}[:,j]$ est la direction de la $j$-ième
feature dans l'espace d'origine (convention Anthropic/Bricken 2023
\cite{bricken2023monosemanticity}). \textbf{SAELens \cite{bloom2024saelens}} (v6) stocke en revanche
\texttt{W\_dec} de forme \texttt{[d\_hidden, d\_model]} : la reconstruction est
alors $\hat{x} = \texttt{sae.W\_dec.T} \mathbin{@} z + \texttt{sae.b\_dec}$.
\end{definitionbox}

\subsection{Architectures des SAE}

\subsubsection{Standard SAE (ReLU, pénalité $L_1$)}
\[
  z = \text{ReLU}(W_\text{enc}(x - b_\text{pre}) + b_\text{enc}), \qquad
  \mathcal{L}_\text{Standard} = \|x - \hat{x}\|_2^2 + \lambda \|z\|_1
\]
Le gradient $\lambda\,\text{sgn}(z)$ induit un biais de rétrécissement
(\emph{shrinkage}) et une vulnérabilité maximale aux latents morts
\cite{cunningham2023sparse}.

\subsubsection{Gated SAE}
Sépare la décision d'activation (\emph{gate}) de sa magnitude
\cite{rajamanoharan2024gated} :
\[
  z = \ind(\pi_\text{gate}(x) > 0) \odot \text{ReLU}(\pi_\text{mag}(x)), \qquad
  \mathcal{L}_\text{Gated} = \|x - \hat{x}\|_2^2 + \lambda \|\text{ReLU}(\pi_\text{gate}(x))\|_1
\]
La pénalité $L_1$ ne porte que sur la porte : les magnitudes actives ne
subissent aucun rétrécissement.

\subsubsection{Top-K et BatchTopK}
Le Top-K contraint exactement $k$ latents actifs par exemple. Le
\textbf{BatchTopK} \cite{gao2024topk, bussmann2025batchtopk}, désormais
recommandé par SAELens comme architecture de référence, relaxe cette
contrainte à l'échelle du batch, avec une perte auxiliaire AuxK pour éviter
les latents morts :
\[
  \mathcal{L}_\text{BatchTopK} = \|x - \hat{x}\|_2^2 + \alpha\,\mathcal{L}_\text{AuxK}(e), \qquad
  \mathcal{L}_\text{AuxK}(e) = \bigl\|e - W_\text{dec}[:,\mathcal{S}]\,\text{TopK}(\pi_{\mathcal{D}_\text{dead}}, k_\text{aux})\bigr\|_2^2
\]

\subsubsection{JumpReLU}
Restaure une parcimonie adaptative sans pénalité $L_1$
\cite{rajamanoharan2024jumprelu} :
\[
  z_i = \pi_i(x) \cdot H(\pi_i(x) - \theta_i)
\]
où $\theta$ est un vecteur de seuils appris et $H$ la fonction de Heaviside.
La discontinuité rend le gradient nul presque partout ; les auteurs
utilisent un estimateur passe-droit (STE) avec approximation par noyau (KDE)
pour mettre à jour $\theta$.

\subsubsection{Matryoshka SAE}
Segmente le dictionnaire en sous-groupes emboîtés
$M = \{m_1 < m_2 < \dots < d_\text{hidden}\}$ et somme les erreurs de
reconstruction à chaque palier \cite{bussmann2025matryoshka} :
\[
  \mathcal{L}_\text{Matryoshka} = \sum_{m \in M} \omega_m \bigl\|x - (W_{\text{dec},\,0:m}\,z_{0:m} + b_\text{dec})\bigr\|_2^2
\]
Cette contrainte force les premières dimensions à encoder l'information la
plus saillante, avec une granularité adaptative à la clé (troncature sans
réentraînement).

\begin{table}[h!]
\centering
\caption{Comparaison analytique des architectures SAE}
\label{tab:sae-archs}
\vspace{0.5em}\small
\begin{tabular}{@{}lllll@{}}
\toprule
\textbf{Architecture} & \textbf{Mécanisme} & \textbf{Pénalité}
  & \textbf{Shrinkage} & \textbf{Statut} \\ \midrule
Standard  & $\text{ReLU}(\pi)$                          & $\lambda\|z\|_1$          & Non  & Baseline \\
Gated     & $\ind(\pi_g>0) \odot \text{ReLU}(\pi_m)$   & $\lambda\|\text{ReLU}(\pi_g)\|_1$ & Oui & Déprécié\\
Top-K     & Troncature $k$ max                          & AuxK                      & Oui  & Déprécié \\
BatchTopK & Troncature $n{\times}k$ max (batch)         & AuxK                      & Oui  & \textbf{Utilisé} \\
JumpReLU  & $\pi \cdot H(\pi - \theta)$                 & $\lambda\|z\|_0$ (STE)    & Oui  & Cœur préentraîné \\
Matryoshka & BatchTopK + nesting                        & $\sum_{m\in M}\|\cdot\|^2$ & Varie & Non implémenté \\ \bottomrule
\end{tabular}
\end{table}

\subsubsection{Le problème des latents morts}

Un latent $j$ est dit \emph{mort} si son activation reste nulle sur un
horizon de batches consécutifs. Pour un Standard SAE, le gradient de la
pénalité $L_1$ pousse monotoniquement les poids vers zéro dès que la
pré-activation devient négative sur tout le batch, ce qui sature
irréversiblement la ReLU. Le diagnostic empirique se fait via
$\rho_\text{dead} = \frac{1}{d_\text{hidden}} \sum_j \ind[\max_x z_j(x) = 0]$ ;
un $\rho_\text{dead} > 5\%$ signale une compression excessive. Chaque
architecture propose son propre mécanisme anti-mort (séparation porte/magnitude
pour Gated, perte AuxK pour Top-K/BatchTopK, \emph{pre-act loss} et
initialisation de seuil pour JumpReLU) --- \emph{cf.}~état de l'art complet en
annexe des sources bibliographiques du stage.

\subsection{Panorama des modèles d'embeddings de texte français}

\subsubsection{CamemBERT 2.0}
Publié par l'équipe ALMAnaCH (Inria) \cite{martin2024camembert}, CamemBERTav2 (base DeBERTaV3, 184M
paramètres) combine un masquage MLM agressif (40\%), un mécanisme
\emph{Disentangled Attention}, et une spécialisation sur des corpus techniques
et scientifiques français (HALvesting). Il constitue l'état de l'art
encodeur pour le français technique en dessous de 200M paramètres.

\subsubsection{Alternatives multilingues (MTEB)}
Le benchmark MTEB-French \cite{kerboua2024mtebfrench} et son extension
multilingue MMTEB \cite{enevoldsen2025mmteb} standardisent l'évaluation sur 8
catégories de tâches. Le panorama dressé en phase préliminaire (Table
\ref{tab:mteb}) couvrait la famille F2LLM-v2 (du 80M au 8B de paramètres,
support natif du \emph{Matryoshka Representation Learning}), Qwen3-Embedding,
multilingual-e5, et CamemBERTav2.

\begin{table}[h!]
\centering
\caption{Panorama des modèles d'embeddings candidats pour le français (relevé bibliographique préliminaire)}
\label{tab:mteb}
\vspace{0.4em}\small
\begin{tabular}{@{}llll@{}}
\toprule
\textbf{Modèle} & \textbf{Params} & \textbf{Points forts} & \textbf{Utilisé dans ce stage ?} \\ \midrule
F2LLM-v2-8B/14B & 8-14B & SOTA global multilingue & Non (contrainte de calcul) \\
Qwen3-Embedding & 0.6-8B & MRL natif, instruction-based & Non \\
CamemBERTav2 & 184M & SOTA $<$200M, français technique & Non \\
\textbf{bge-m3} & 560M & Multilingue, pooling [CLS] & \textbf{Oui} (Pipeline 2) \\
\textbf{F2LLM-v2-80/160/330M} & 80-330M & MRL, léger, déployable & \textbf{Oui} (Pipeline 2) \\
\bottomrule
\end{tabular}
\end{table}

\begin{diagnosticbox}[Écart entre relevé bibliographique et choix effectif]
Le relevé bibliographique préliminaire du stage plaçait CamemBERTav2 comme
candidat privilégié pour un corpus français technique. En pratique,
l'architecture retenue s'appuie sur \textbf{GemmaScope} (SAE préentraîné sur
Gemma-3, disponible \emph{prêt à l'emploi} avec un catalogue de labels
Neuronpedia déjà constitué) pour le pipeline principal, et sur
\textbf{F2LLM-v2}/\textbf{bge-m3} pour le second pipeline --- CamemBERTav2 n'a
finalement pas été intégré, faute d'un SAE préentraîné équivalent disponible
pour ce backbone (l'intégralité de l'écosystème GemmaScope/Neuronpedia, qui
motive le choix de Gemma-3, n'existe pas pour CamemBERT). Ce choix, motivé par
la disponibilité d'infrastructure plutôt que par un score MTEB, est assumé et
documenté comme limite au chapitre~\ref{sec:discussion}.
\end{diagnosticbox}

\subsubsection{Matryoshka Representation Learning et architectures décodeurs}
Le MRL \cite{kusupati2022mrl} entraîne un modèle d'embeddings de sorte que des
sous-vecteurs de dimensions imbriquées soient tous des points d'ancrage
sémantiques valides, via une perte contrastive multi-résolution
$\mathcal{L}_{MRL} = \sum_{d \in \mathcal{D}} \omega_d \cdot \mathcal{L}_\text{contrastive}(z_{1:d})$.
F2LLM-v2 en bénéficie nativement, ce qui a directement motivé le test empirique
de troncature d'embedding présenté au chapitre~\ref{sec:resultats}.

\subsection{Le SAE résiduel à cœur gelé : deux généalogies concurrentes}
\label{subsec:sae-boost-litt}

La piste architecturale retenue pour ce stage --- geler un SAE préentraîné et
entraîner une extension sur son résidu de reconstruction --- a été développée
en parallèle par ce projet et, dans la littérature publiée, par
\textbf{SAE Boost} \cite{koriagin2025saeboost} (Koriagin et al., juillet 2025) :

\begin{align*}
  e &= x - \hat{x}_\text{core} & \text{(résidu du SAE gelé)} \\
  z_\text{res} &= \text{BatchTopK}_{k_\text{res}}\bigl(\text{ReLU}(W_\text{enc}^{(\text{res})} x + b_\text{enc}^{(\text{res})})\bigr) \\
  \mathcal{L}_\text{Boost} &= \|e - \hat{e}\|_2^2 / \Var(e)
\end{align*}

Koriagin et al. rapportent des gains de +25\% à +59\% en variance expliquée
sur des domaines cibles (chimie, russe, débats ONU) pour moins de 1\% de
dégradation sur le domaine général. Cette architecture résout l'interférence
de routage \emph{par construction} : aucun budget de sparsité n'est partagé
entre cœur et extension. C'est l'architecture concrètement implémentée et
validée à l'échelle dans ce stage (\emph{cf.}~\S\ref{subsec:frozencore}) ---
et non une architecture Matryoshka pleinement hiérarchique et emboîtée, la
piste initialement envisagée en phase bibliographique préliminaire mais dont
la complexité d'entraînement et le besoin en données n'étaient pas justifiés
par les résultats obtenus (\emph{cf.}~chapitre~\ref{sec:resultats}).

\subsection{Sanity checks : les métriques proxy suffisent-elles ?}

Korznikov et al. \cite{korznikov2025sanitychecks} montrent qu'un SAE dont le
décodeur est figé à une initialisation aléatoire (jamais entraîné) atteint,
sur les métriques proxy usuelles (variance expliquée, AutoInterp, sparse
probing), des performances étonnamment proches d'un SAE pleinement entraîné.
Ce résultat impose un critère de validation supplémentaire : toute
amélioration mesurée doit être évaluée \emph{relativement} à cette baseline
aléatoire gelée, et non uniquement en absolu. Ce protocole a été directement
reproduit dans ce stage (\emph{cf.}~\S\ref{subsec:sanity-check-resultats}).

\subsection{Interprétabilité au niveau documentaire}

L'application des SAE à des embeddings de documents entiers (plutôt qu'à des
activations token-level) constitue une extension récente de la littérature
\cite{jiang2024sae_embeddings, decodedense2025}. Le pipeline documentaire
distingue \textbf{interprétabilité globale} (que représente la feature $j$ sur
l'ensemble du corpus ? --- mesurée par juge LLM, protocole \emph{odd-one-out})
et \textbf{interprétabilité locale} (quelles features sont actives pour un
document donné ? --- mesurée par max-pooling documentaire des activations
token-level). Le protocole \emph{odd-one-out}
\cite{templeton2024scaling} présente au juge 9 extraits activant fortement une
même feature et 1 intrus aléatoire ; le juge doit identifier l'intrus, ce qui
ne peut réussir, par construction, que si les 9 extraits partagent un concept
identifiable.

\newpage
% ─────────────────────────────────────────────────────────────────────────────
\section{Contribution : architecture et méthodologie}
\label{sec:contribution}
% ─────────────────────────────────────────────────────────────────────────────

\subsection{Vue d'ensemble du système}

Le système implémente deux pipelines d'analyse interprétable de mails
clients EDF, partageant la même infrastructure de corpus, de stockage et
d'évaluation.

\begin{figure}[h!]
\centering
\begin{tikzpicture}[
  node distance=0.7cm and 0.8cm,
  box/.style={draw, rounded corners=2pt, minimum width=2.6cm, minimum height=0.8cm,
    align=center, font=\scriptsize, fill=gray!10},
  arrow/.style={-Stealth, thick}
]
  \node[box] (corpus) {Corpus principal\\Mails originaux +\\variantes augmentées};
  \node[box, below left=1.2cm and -0.3cm of corpus] (p1) {Pipeline 1\\Gemma-3-12B-it\\+ GemmaScope};
  \node[box, below right=1.2cm and -0.3cm of corpus] (p2) {Pipeline 2\\F2LLM-v2 / bge-m3\\+ PhraseLevelSAE};
  \node[box, below=0.9cm of p1] (ext) {Extension\\FrozenCoreResidualSAE\\(entraînée sur résidu)};
  \node[box, below=0.9cm of p2] (train) {SAE entraîné\\from-scratch};
  \node[box, below=0.9cm of ext] (pool1) {Max-pool\\documentaire};
  \node[box, below=0.9cm of train] (pool2) {Max-pool\\documentaire};
  \node[box, below=0.9cm of pool1] (lab1) {Labels : Neuronpedia\\+ juge LLM local};
  \node[box, below=0.9cm of pool2] (lab2) {Labels : juge LLM\\local (odd-one-out)};

  \draw[arrow] (corpus) -- (p1);
  \draw[arrow] (corpus) -- (p2);
  \draw[arrow] (p1) -- (ext);
  \draw[arrow] (p2) -- (train);
  \draw[arrow] (ext) -- (pool1);
  \draw[arrow] (train) -- (pool2);
  \draw[arrow] (pool1) -- (lab1);
  \draw[arrow] (pool2) -- (lab2);
\end{tikzpicture}
\caption{Architecture générale des deux pipelines d'analyse}
\label{fig:pipeline}
\end{figure}

\subsection{Architecture \texorpdfstring{$\textit{FrozenCoreResidualSAE}$}{FrozenCoreResidualSAE}}
\label{subsec:frozencore}

Pour aligner les embeddings d'origine (conçus majoritairement sur des corpus
anglophones) avec les spécificités sémantiques et techniques du corpus cible,
sans induire d'oubli catastrophique sur les capacités générales du SAE
préentraîné \cite{lieberum2024gemmascope}, le Pipeline~1 gèle le dictionnaire du cœur GemmaScope
($W_\text{dec}^\text{core}$) et instancie une extension de $D_\text{extra}$
features dédiées, entraînée exclusivement sur le résidu non reconstruit :
\[
  e = x - \hat{x}_\text{core} = x - \Bigl(b_\text{dec}^\text{core} +
  \sum_{j \in \mathcal{A}_\text{core}} z_j^\text{core}\, w_{\text{dec},j}^\text{core}\Bigr)
\]
\[
  z_\text{extra} = \text{BatchTopK}_{K_\text{extra}}\bigl(\text{ReLU}(W_\text{enc}^\text{extra} x + b_\text{enc}^\text{extra})\bigr), \qquad
  \hat{x} = \hat{x}_\text{core} + W_\text{dec}^\text{extra}\, z_\text{extra}
\]
Les poids $W_\text{dec}^\text{extra}$ sont initialisés par PCA sur un
sous-ensemble de calibration en français ; l'absence de biais additionnel
($b_\text{dec}^\text{extra}=0$) garantit que sur un flux hors-domaine
l'extension reste quasi-inactive, préservant la géométrie de base. Cette
architecture correspond exactement à celle décrite pour SAE Boost
(\S\ref{subsec:sae-boost-litt}) : les résultats de ce stage constituent donc,
\emph{de facto}, une réplication indépendante et une extension empirique de
cette architecture sur un nouveau domaine (emails clients français) et à une
échelle de modèle non couverte par la publication d'origine
(1B/4B/12B~paramètres, \emph{cf.}~\S\ref{subsec:model-scale}).

\subsection{Métriques d'évaluation}

\subsubsection{Fidélité topologique locale : \texorpdfstring{$\rho_\text{SAE}$}{rho\_SAE}}
Pour vérifier que la projection SAE ne dénature pas la géométrie de l'espace
d'origine, on évalue la corrélation de rang de Spearman entre les matrices de
similarité cosinus calculées avant et après passage par le SAE, sur un
sous-ensemble de $N$ documents échantillonnés :
\[
  \rho_\text{SAE} = 1 - \frac{6\sum_{i=1}^n d_i^2}{n(n^2-1)}, \qquad
  d_i = \text{rang}(M_\text{raw}^{(i)}) - \text{rang}(M_\text{sae}^{(i)}), \qquad
  n = \binom{N}{2}
\]
Un $\rho_\text{SAE} \to 1$ valide empiriquement la conservation de la
topologie de l'espace de base.

\subsubsection{Sonde logistique (probing)}
Pour vérifier que la décomposition parcimonieuse préserve le pouvoir
discriminant, deux régressions logistiques multinomiales sont ajustées en
parallèle : $\text{Acc}_\text{raw}$ sur les embeddings denses d'origine
(borne supérieure oracle) et $\text{Acc}_\text{sae}$ sur les codes SAE. Le
différentiel $\Delta = \text{Acc}_\text{raw} - \text{Acc}_\text{sae}$ mesure
la perte d'information discriminante induite par la parcimonie. Au-delà de la
conservation de performance, cette sonde rend le modèle intrinsèquement
explicable : les coefficients $\beta_j^{(k)}$ les plus élevés isolent
directement les concepts sémantiques qui dictent la décision.

\subsubsection{Auto-interprétation par juge LLM}
Le taux d'interprétabilité rapporté dans tout ce rapport est mesuré par le
protocole \emph{odd-one-out} : pour chaque feature échantillonnée, un juge LLM
(le même modèle que l'extracteur, sauf mention contraire) reçoit 9 documents
activant fortement la feature et 1 document-contrôle (activation quasi
nulle), et doit désigner l'intrus. Un juge qui réussit démontre l'existence
d'un concept partagé identifiable ; l'échec ne prouve pas nécessairement
l'absence de concept (bruit du protocole, \emph{cf.}~\S\ref{subsec:robustesse}).

\subsection{Construction du corpus}

Le corpus principal combine les mails « originaux » (3\,480 documents,
\emph{cf.}~encadré \S\ref{sec:etat-art}) et leurs variantes augmentées
(39\,949 variantes acceptées sur 45\,240 générées, soit un taux de rejet de
qualité de 11,7\% --- \emph{cf.}~\S\ref{subsec:augmentation}), produites selon
13 axes de perturbation contrôlée (émotion, registre, orthographe, urgence).
Le split train/test est effectué au niveau du \textbf{mail d'origine} (et non
de la ligne) pour éviter qu'une variante augmentée d'un mail de test ne fuite
dans le train --- une variante et son mail source partagent un contenu
factuel quasi identique, ce qui biaiserait artificiellement les métriques de
classification en aval si elles se retrouvaient de part et d'autre du split.

\subsubsection{Garde-fou qualité de la génération augmentée}
\label{subsec:augmentation}
Chaque variante générée est soumise à un contrôle automatique (longueur
minimale, ratio de longueur par rapport au mail source dans $[0{,}4\,;\,2{,}5]$,
non-identité stricte, préservation des entités numériques du mail d'origine
sauf pour l'axe orthographe). 11,7\% des générations sont rejetées par ce
garde-fou, sans impact sur les résultats rapportés : les variantes rejetées
ne sont jamais chargées en aval de la génération.

\subsection{Protocole de validation et infrastructure}

L'ensemble des expériences a été exécuté sur un cluster SLURM à plusieurs
partitions GPU (A100/H100), sans accès réseau direct depuis les nœuds de
calcul --- toutes les dépendances (modèles, données) doivent être
prépositionnées sur disque avant soumission. Gemma-3 présente des activations
de norme très élevée (\emph{massive activations}) dans son flux résiduel ; une
précision fp16 (plage dynamique $\sim 65\,504$) provoque un dépassement
silencieux vers \texttt{inf}/\texttt{nan} sur ces valeurs, corrompant tout
entraînement en aval sans erreur explicite --- le projet utilise bf16 par
défaut partout, qui partage la plage d'exposant de fp32.

Le taux d'interprétabilité étant une proportion mesurée sur un échantillon de
$n=150$ features, chaque comparaison entre deux configurations est testée
statistiquement (test $z$ à deux proportions indépendantes par défaut ; test de
McNemar pour les comparaisons appariées sur les mêmes features ;
Cochran-Armitage pour les plans à niveaux ordonnés --- \emph{cf.}~audit
méthodologique \S\ref{subsec:audit-stats}), avec un seuil conventionnel de
signification $|z| > 1{,}96$.

\newpage
% ─────────────────────────────────────────────────────────────────────────────
\section{Résultats expérimentaux}
\label{sec:resultats}
% ─────────────────────────────────────────────────────────────────────────────

\subsection{Diagnostic et correction du problème initial}
\label{subsec:diagnostic-initial}

Le pipeline complet fonctionnait de bout en bout, mais le run disponible en
début de cette phase (10 features jugées, corpus generic
énergie/sport/support) ne passait le test \emph{odd-one-out} que dans
\textbf{20,0\% des cas} (2/10). L'inspection qualitative des exemples présentés
au juge a révélé l'absence de concept commun identifiable entre les 9 extraits
positifs d'une même feature (\emph{ex.} : rappel produit iPad, système
carcéral norvégien, recette de cuisine et agriculture canadienne présentés
ensemble). La lecture du code d'assemblage du corpus a montré la cause
racine : le réservoir de résidus servant à entraîner l'extension était
constitué \textbf{exclusivement} de texte générique (Wikipédia filtré par
mots-clés) --- les mails du domaine cible n'entraient jamais dans les données
d'entraînement de l'extension, seulement dans une visualisation post-hoc.

\begin{resultbox}[Correction et validation]
Une fois le corpus principal reconstruit (mails + variantes augmentées,
\S\ref{sec:contribution}) et le nombre de features jugées porté à 150 pour une
puissance statistique correcte, le taux d'interprétabilité passe à
\textbf{45,3\% (68/150)} --- plus du double du taux initial, à volume de
tokens comparable. Faire varier ce volume d'un facteur 20 (100k à 2M tokens),
à domaine fixé, ne produit \textbf{aucun effet mesurable} (40,7\% / 45,3\% /
44,7\%, écarts non significatifs). \textbf{Le problème initial était donc un
défaut de domaine du corpus, pas un déficit de volume d'entraînement.}
\end{resultbox}

\subsection{Synthèse des ablations d'hyperparamètres du SAE}
\label{subsec:synthese-ablations}

Une fois le domaine du corpus corrigé, une campagne d'ablations systématique
a testé, pour chaque configuration, une modification \textbf{isolée} par
rapport au run principal (référence : 45,3\%, 68/150), toutes comparées par
test $z$ à deux proportions :

\begin{table}[h!]
\centering
\caption{Synthèse des 14 ablations d'hyperparamètres --- taux d'interprétabilité (n=150)}
\label{tab:synthese}
\vspace{0.4em}\small
\begin{tabular}{@{}lccc@{}}
\toprule
\textbf{Configuration} & \textbf{Taux} & \textbf{z vs référence} & \textbf{Signif.} \\ \midrule
Volume 100k tokens & 40,7\% & 0,82 & non \\
Volume 2M tokens & 44,7\% & 0,12 & non \\
Volume 25M tokens (+filler FineWeb2-fr) & 54,0\% & --1,50 & non \\
Largeur SAE core 65k & 43,3\% & 0,35 & non \\
Largeur SAE core 262k & 46,7\% & --0,23 & non \\
Époques $\times$4 (40/100) & 41,3\% & 0,70 & non \\
Capacité combinée ($D{=}2048$, $K{=}64$) & 40,0\% & 0,93 & non \\
$D_\text{extra}=2048$ seul ($K$ fixe) & 46,0\% & --0,12 & non \\
$K_\text{extra}=5$ (optimum du papier SAE Boost) & 54,7\% & --1,62 & non \\
Graine d'entraînement (seed 123) & 47,3\% & --0,35 & non \\
Combiné v12 (65k+époques+capacité, tranche 1-150) & 44,0\% & 0,23 & non \\
\textbf{Sanity check : décodeur figé aléatoire} & \textbf{29,3\%} & \textbf{2,86} & \textbf{oui} \\
\textbf{Échelle du modèle --- 4B} & \textbf{28,0\%} & \textbf{3,12} & \textbf{oui} \\
\textbf{Échelle du modèle --- 1B} & \textbf{12,0\%} & \textbf{6,38} & \textbf{oui} \\
\bottomrule
\end{tabular}
\end{table}

La graine d'entraînement, en particulier, était motivée par la littérature sur l'instabilité des features individuelles d'un SAE \cite{unstablefeatures2026, identifiablesae2026}.

\textbf{Constat central} : sur les 14 configurations testées, \textbf{11
n'ont aucun effet significatif} sur le taux d'interprétabilité --- ni la
largeur du SAE core, ni le volume de tokens, ni la capacité ou la parcimonie
de l'extension (isolées ou combinées), ni la graine d'entraînement, ne
déplacent la mesure au-delà du bruit d'échantillonnage attendu à $n=150$. Les
3 configurations significatives se distinguent nettement des autres, tant en
magnitude qu'en nature : un sanity check architectural et un unique
hyperparamètre, l'\textbf{échelle du modèle}, discuté ci-dessous.

\subsection{Sanity check : le SAE bat-il un décodeur figé aléatoire ?}
\label{subsec:sanity-check-resultats}

Reproduisant le protocole de Korznikov et al.~\cite{korznikov2025sanitychecks}
(\S\ref{sec:etat-art}), un décodeur $W_\text{dec}^\text{extra}$ figé à une
initialisation aléatoire (jamais entraîné) atteint tout de même
\textbf{29,3\%} d'interprétabilité (44/150) --- nettement en dessous du SAE
pleinement entraîné (45,3\%, écart significatif, $z=2,86$) mais loin d'être
nul. Sur la métrique de classification (séparabilité des 14 classes d'axes
de perturbation), l'écart se resserre fortement (93,5\% entraîné contre
91,2\% figé) : \textbf{la classification en aval résiste beaucoup mieux à un
décodeur aléatoire que l'interprétation qualitative}, répliquant partiellement
le constat du papier sur des métriques proxy insuffisantes.

\subsection{L'échelle du modèle : le levier déterminant}
\label{subsec:model-scale}

Trois échelles du modèle extracteur/juge ont été testées, chacune avec son
propre GemmaScope dédié (gemma-scope-2-1b-it, -4b-it, -12b-it), toutes choses
égales par ailleurs :

\begin{table}[h!]
\centering
\caption{Effet dose-réponse de l'échelle du modèle}
\label{tab:model-scale}
\vspace{0.4em}\small
\begin{tabular}{@{}lccccc@{}}
\toprule
\textbf{Modèle} & \textbf{$d_\text{model}$} & \textbf{Taux interp.} & \textbf{$z$ vs 12B} & \textbf{Acc. classification} \\ \midrule
gemma-3-1b-it & 1152 & 12,0\% (18/150) & 6,38 & 88,2\% \\
gemma-3-4b-it & 2560 & 28,0\% (42/150) & 3,12 & 92,0\% \\
gemma-3-12b-it (référence) & 3840 & 45,3\% (68/150) & --- & 93,5\% \\
\bottomrule
\end{tabular}
\end{table}

\begin{resultbox}[Résultat central du stage]
La progression 12,0\% $\to$ 28,0\% $\to$ 45,3\% est \textbf{monotone et
significative à chaque palier}, y compris la comparaison directe 1B contre 4B
($z=-3{,}46$). Un test de tendance de Cochran-Armitage --- statistique unique,
mieux adaptée qu'une série de comparaisons par paires à ce plan à niveaux
ordonnés --- confirme un effet extrêmement robuste :
$z \approx 6{,}40$, $p \approx 1{,}6\times10^{-10}$ (quasiment identique avec
un barème linéaire ou log-paramètres). C'est, de très loin, le plus fort
effet mesuré dans tout ce stage.
\end{resultbox}

La séparabilité linéaire (classification, indépendante du juge LLM) suit une
pente beaucoup plus douce (88,2\%~$\to$~93,5\%, 5,3~points d'écart total,
contre 33,3~points pour l'interprétabilité qualitative) : les représentations
latentes restent largement exploitables même à 1B, alors que le jugement
qualitatif s'effondre. Cette dissociation, ainsi que l'inspection qualitative
des hypothèses de diffing générées (cohérentes à 12B, confuses à 1B ---
inversion logique cause/conséquence observée), suggère que c'est la
\textbf{capacité de raisonnement du juge} --- formuler et vérifier un concept
partagé entre 9 exemples --- qui limite le score à petite échelle, plus que
la qualité intrinsèque des features extraites.

\subsection{Robustesse du protocole de jugement}
\label{subsec:robustesse}

Deux tests évaluent la sensibilité du protocole \emph{odd-one-out} à des
perturbations de surface qui ne devraient, en principe, pas affecter le
jugement :

\begin{itemize}
  \item \textbf{Réordonnancement} : en répétant la question 5 fois par
    feature avec un ordre de présentation différent, seules 30,7\% des
    features obtiennent une décision unanime. Le taux agrégé bouge peu
    (single-shot 45,3\% contre vote majoritaire 48,7\%, test de McNemar
    $p=0,560$, non significatif) mais \textbf{31,3\% des features changent
    individuellement de statut} selon l'ordre de présentation.
  \item \textbf{Langue du juge} \cite{resck2025multilingual, multilingualconcepts2025} : traduire les exemples et le prompt en
    anglais ne change pas significativement le taux agrégé (français 46,9\%
    contre anglais 45,5\%, McNemar $p=0,894$), mais \textbf{38,6\% des
    features changent de statut} selon la langue --- un taux de bruit
    supérieur à celui du simple réordonnancement.
\end{itemize}

Ces deux résultats convergent vers la même conclusion : le protocole
\emph{odd-one-out} à décision unique est \textbf{générique\-ment bruyant} face
à toute perturbation de surface (ordre, langue), pas spécifiquement biaisé
dans une direction donnée. Un vote majoritaire sur plusieurs répétitions
devrait être préféré à une décision \emph{greedy} unique en usage de
production.

\subsection{Qualité de l'explication document-level}

Deux protocoles complémentaires valident que les features citées comme
explication d'une décision de classification y jouent un rôle causal réel,
et non de simples labels plausibles sans lien :

\begin{itemize}
  \item \textbf{Fidélité par ablation} : mettre à zéro les $K$ features les
    plus contributives à une décision de classification produit une chute de
    probabilité prédite nettement supérieure à l'ablation de $K$ features
    aléatoires ou des $K$ moins contributives --- confirmant que les features
    citées comme explication pilotent réellement la décision.
  \item \textbf{Plausibilité par choix forcé} : un juge LLM, confronté aux
    vraies features actives d'un document et à un jeu de features-leurres,
    préfère significativement les vraies features, plus souvent que le
    hasard.
\end{itemize}

\subsection{Fidélité du steering (\texttt{steer\_and\_decode})}

Le steering (amplifier ou supprimer une feature puis décoder le résultat)
n'avait jamais été testé au-delà d'une vérification géométrique superficielle.
Un test dédié, suppriment les 10 features les plus explicatives d'une
intention puis décodant/ré-encodant le résultat, révèle un comportement
\textbf{très hétérogène selon l'intention} : le \emph{round-trip}
decode/ré-encode neutralise quasiment l'intervention pour 2 intentions sur 4
testées (ratio 0,00--0,02$\times$ par rapport à l'ablation directe en espace de
code), la préserve pour une troisième (0,90$\times$), et l'amplifie pour la
dernière (1,74$\times$). Le steering décodé n'est donc \textbf{pas} un
mécanisme d'intervention causale fiable et prévisible à partir du simple test
d'ablation en espace de code.

\subsection{Retrieval quantitatif (Latent Terms)}

Une évaluation quantitative --- jamais réalisée jusqu'ici --- du retrieval
BM25 sur vocabulaire latent SAE \cite{clavie2026latentterms} montre une
\textbf{précision parfaite (1,00)} sur 3 intentions testées sur 4, nettement
supérieure à une baseline TF-IDF classique sur des requêtes reformulées (pas
de reprise littérale des mots-clés) --- généralisation sémantique réelle. Un
échec complet sur la quatrième intention (0,00) a été diagnostiqué
précisément : un seul document du corpus partage une intersection non nulle
avec les features latentes activées par la requête --- une limite
structurelle du BM25 sur vocabulaire très parcimonieux, pas un défaut de
compréhension sémantique.

\subsection{Choix du backbone d'embedding (Pipeline 2)}

Quatre backbones ont été comparés pour le Pipeline~2 : F2LLM-v2 (80M, 160M,
330M) et bge-m3. \textbf{bge-m3 domine nettement} sur la plupart des
métriques (NMSE $-18,8\%$ par rapport au meilleur F2LLM, taux de features
mortes 4 à 5 fois plus faible, meilleure séparabilité du corpus de diffing)
--- un candidat par défaut recommandé pour une suite du projet. Ce résultat a
nécessité un correctif de code : le pooling documentaire était câblé en dur
sur le dernier token (adapté à un backbone décodeur causal comme F2LLM),
incorrect pour un encodeur bidirectionnel comme bge-m3 (pooling \texttt{[CLS]}
attendu).

\subsection{Balayage de la dimension d'embedding (MRL)}

La dimension d'embedding (\texttt{MATRYOSHKA\_DIM}, F2LLM) n'avait jamais été
variée --- un fait notable découvert en creusant la question : pour
F2LLM-80M, la dimension par défaut égale exactement la dimension complète du
modèle, rendant toute \og{}troncature\fg{} antérieure un artefact sans effet
réel. Un balayage sur F2LLM-160M (dimensions 64, 128, 320, 640) montre une
\textbf{dégradation graduelle, pas abrupte} : 10\% des dimensions (64/640)
conservent déjà 94\% de la performance de classification du plein embedding
--- cohérent avec (sans le démontrer formellement, faute de confirmation d'un
entraînement MRL strict côté F2LLM) un comportement de type Matryoshka.

\subsection{Audit de la méthodologie statistique}
\label{subsec:audit-stats}

Un audit rétroactif de la rigueur statistique des comparaisons de ce stage a
identifié et corrigé trois lacunes :

\begin{enumerate}
  \item Les comparaisons appariées (biais multilingue, robustesse du juge)
    avaient été analysées par un test à deux proportions indépendantes plutôt
    que par le test de \textbf{McNemar}, statistiquement approprié lorsque
    les deux conditions portent sur les mêmes 150 features. Le recalcul ne
    change aucune conclusion ($p=0{,}894$ et $p=0{,}560$) mais corrige la
    méthode.
  \item Aucune \textbf{correction pour comparaisons multiples} n'avait été
    appliquée aux $\sim$15 tests d'ablation de ce chapitre (contrairement au
    diffing par feature, qui utilise déjà Benjamini-Hochberg). Sans
    conséquence pratique ici (le seul résultat significatif l'est à
    $p<10^{-9}$), mais une lacune à corriger pour toute extension future.
  \item L'effet dose-réponse de l'échelle du modèle a été re-confirmé par un
    \textbf{test de tendance} (Cochran-Armitage) plutôt qu'une série de tests
    par paires --- statistique unique, plus puissante et mieux adaptée à un
    plan à niveaux ordonnés.
\end{enumerate}

\newpage
% ─────────────────────────────────────────────────────────────────────────────
\section{Discussion et conclusion}
\label{sec:discussion}
% ─────────────────────────────────────────────────────────────────────────────

\subsection{Bilan par rapport aux objectifs initiaux}

\begin{table}[h!]
\centering
\caption{Bilan par rapport aux étapes de l'offre de stage}
\vspace{0.4em}\small
\begin{tabular}{@{}p{5.5cm}p{8.5cm}@{}}
\toprule
\textbf{Étape prévue} & \textbf{Réalisation} \\ \midrule
Partie bibliographique sur les SAE & Faite --- élargie à un panorama des
  architectures, des embeddings, des lois d'échelle, et de la littérature de
  sanity-check récente (chapitre~\ref{sec:etat-art}). \\
Choix d'une implémentation existante & GemmaScope + SAELens comme cœur
  préentraîné, complété d'une extension développée pendant le stage. \\
Conversion des données en embeddings & Deux pipelines indépendants
  (Gemma-3/GemmaScope, F2LLM/bge-m3). \\
Entraînement d'un SAE & Fait, avec diagnostic et correction d'un défaut de
  corpus majeur en cours de route. \\
Analyses, interprétation, visualisation & Dashboard interactif, juge LLM
  automatisé, retrieval, diffing, sondes de classification. \\
Rédaction d'un rapport & Ce document. \\
\bottomrule
\end{tabular}
\end{table}

\subsection{Points résolus}

Le défaut de conception du corpus initial a été diagnostiqué et corrigé, avec
un gain d'interprétabilité mesuré et statistiquement établi. L'architecture à
cœur gelé a été validée à l'échelle, avec une reconnaissance explicite de son
équivalence structurelle à SAE Boost. Une campagne d'ablations
particulièrement large a permis d'établir, avec une rigueur statistique
correcte (McNemar, correction de multiplicité identifiée, test de tendance),
qu'aucun hyperparamètre du SAE ne module significativement le résultat une
fois le domaine corrigé --- et qu'un unique levier, l'échelle du modèle,
domine tous les autres effets mesurés. Les protocoles de fidélité,
plausibilité et robustesse de l'explication ont tous été mis en œuvre et
donnent des résultats positifs ou, à défaut, honnêtement nuancés (steering).

\subsection{Points non résolus et perspectives}

\begin{itemize}
  \item L'architecture Matryoshka pleinement hiérarchique (dictionnaires
    emboîtés à plusieurs niveaux), envisagée en phase bibliographique,
    n'a pas été implémentée --- l'architecture à cœur gelé effectivement
    testée n'a qu'un seul niveau d'extension. Au vu du résultat central de ce
    stage (l'échelle du modèle domine tout effet architectural du SAE), une
    architecture Matryoshka multi-niveaux ne semble pas prioritaire pour la
    suite : l'effort gagnerait probablement davantage à séparer les rôles
    d'extracteur et de juge (utiliser un modèle léger pour l'extraction mais
    un juge plus capable) pour isoler laquelle des deux capacités domine
    réellement l'effet d'échelle observé.
  \item CamemBERTav2 et Qwen3-Embedding, identifiés en phase bibliographique
    comme candidats sérieux, n'ont pas été intégrés faute d'écosystème SAE
    préentraîné équivalent à GemmaScope --- une comparaison chiffrée directe
    reste à faire si l'équipe souhaite objectiver ce choix par rapport au
    score MTEB plutôt que par la disponibilité d'infrastructure.
  \item Le pipeline micro/macro dual-SAE avec alignement AlignedUMAP, la
    modélisation par graphe biparti, et le protocole d'annotation humaine
    (Fleiss $\kappa$) envisagés en phase préliminaire restent des pistes non
    explorées, faute de temps.
  \item Aucune comparaison chiffrée directe n'a été menée contre les
    baselines alternatives de SAE Boost (Extended SAE à initialisation
    aléatoire/most-active, SAE Stitching, fine-tuning complet).
\end{itemize}

\subsection{Conclusion générale}

Ce stage a permis de construire, diagnostiquer et valider empiriquement, sur
un cas d'usage industriel réel (mails clients EDF), une architecture SAE à
cœur gelé pour l'interprétabilité de grands modèles de langage. Le résultat
le plus significatif --- l'effet dominant de l'échelle du modèle
extracteur/juge sur le taux d'interprétabilité mesuré, loin devant tout
hyperparamètre du SAE lui-même --- reformule la question de recherche
initiale : au-delà du choix de l'architecture SAE, c'est la capacité du
modèle de langage sous-jacent à raisonner sur un concept partagé qui
conditionne la qualité de l'interprétation automatique obtenue. Cette
conclusion, obtenue par une démarche d'ablation rigoureuse et statistiquement
auditée, constitue une contribution méthodologique transférable à tout projet
similaire d'interprétabilité appliquée.

\newpage
% =============================================================================
% BIBLIOGRAPHIE
% =============================================================================
\addcontentsline{toc}{section}{Bibliographie}
\begin{thebibliography}{99}\small

\bibitem{elhage2022superposition}
  Elhage, N. et al.
  \textit{Toy Models of Superposition}.
  Transformer Circuits Thread, Anthropic, 2022.

\bibitem{bricken2023monosemanticity}
  Bricken, T. et al.
  \textit{Towards Monosemanticity: Decomposing Language Models With Dictionary Learning}.
  Transformer Circuits Thread, Anthropic, 2023.

\bibitem{cunningham2023sparse}
  Cunningham, H., Ewart, A., Riggs, L., Huben, R., \& Sharkey, L.
  \textit{Sparse Autoencoders Find Highly Interpretable Features in Language Models}.
  ICLR 2024. arXiv:2309.08600.

\bibitem{rajamanoharan2024gated}
  Rajamanoharan, S. et al.
  \textit{Improving Dictionary Learning with Gated Sparse Autoencoders}.
  arXiv:2404.16014, 2024.

\bibitem{gao2024topk}
  Gao, L. et al.
  \textit{Scaling and Evaluating Sparse Autoencoders}.
  arXiv:2406.04093, OpenAI, 2024.

\bibitem{bussmann2025batchtopk}
  Bussmann, B., Nabeshima, N., Karvonen, A., \& Nanda, N.
  \textit{Learning Multi-Level Features with Matryoshka Sparse Autoencoders}.
  ICML 2025. arXiv:2503.17547.

\bibitem{rajamanoharan2024jumprelu}
  Rajamanoharan, S. et al.
  \textit{Jumping Ahead: Improving Reconstruction Fidelity with JumpReLU Sparse Autoencoders}.
  arXiv:2407.14435, DeepMind, 2024.

\bibitem{bussmann2025matryoshka}
  Bussmann, B., Nabeshima, N., Karvonen, A., \& Nanda, N.
  \textit{Learning Multi-Level Features with Matryoshka Sparse Autoencoders}.
  ICML 2025. arXiv:2503.17547.

\bibitem{templeton2024scaling}
  Templeton, A. et al.
  \textit{Scaling Monosemanticity: Extracting Interpretable Features from Claude 3 Sonnet}.
  Anthropic, 2024.

\bibitem{bloom2024saelens}
  Bloom, J., Tigges, C., Duong, A., \& Chanin, D.
  \textit{SAELens}. \url{https://github.com/decoderesearch/SAELens}, 2024.

\bibitem{jiang2024sae_embeddings}
  (Collectif MATS).
  \textit{Interpretable Embeddings with Sparse Autoencoders: A Data Analysis Toolkit}.
  arXiv:2512.10092, 2024.

\bibitem{decodedense2025}
  (Anonyme).
  \textit{Decoding Dense Embeddings: Sparse Autoencoders for Interpreting and Discretizing Dense Retrieval}.
  arXiv:2506.00041, 2025.

\bibitem{martin2024camembert}
  Martin, L. et al.
  \textit{CamemBERT 2.0: A Smarter French Language Model Aged to Perfection}.
  arXiv:2411.08868, Inria/ALMAnaCH, 2024.

\bibitem{kerboua2024mtebfrench}
  Kerboua, I. et al.
  \textit{MTEB-French: Resources for French Sentence Embedding Evaluation and Analysis}.
  arXiv:2405.20468, 2024.

\bibitem{enevoldsen2025mmteb}
  Enevoldsen, K. et al.
  \textit{MMTEB: Massive Multilingual Text Embedding Benchmark}.
  ICLR 2025. arXiv:2502.13595.

\bibitem{kusupati2022mrl}
  Kusupati, A. et al.
  \textit{Matryoshka Representation Learning}.
  NeurIPS 2022.

\bibitem{koriagin2025saeboost}
  Koriagin, N., Aksenov, Y., Laptev, D., Gerasimov, G., Balagansky, N., \& Gavrilov, D.
  \textit{Teach Old SAEs New Domain Tricks with Boosting}.
  arXiv:2507.12990, 2025.

\bibitem{korznikov2025sanitychecks}
  Korznikov, A., Galichin, A., Dontsov, A., Rogov, O., Oseledets, I., \& Tutubalina, E.
  \textit{Sanity Checks for Sparse Autoencoders: Do SAEs Beat Random Baselines?}
  arXiv:2602.14111, 2025.

\bibitem{lieberum2024gemmascope}
  Lieberum, T. et al.
  \textit{Gemma Scope: Open Sparse Autoencoders Everywhere All At Once on Gemma 2}.
  arXiv:2408.05147, Google DeepMind, 2024.

\bibitem{resck2025multilingual}
  Resck, L., Augenstein, I., \& Korhonen, A.
  \textit{Explainability and Interpretability of Multilingual LLMs: A Survey}.
  EMNLP 2025.

\bibitem{clavie2026latentterms}
  Clavié, B. et al.
  \textit{Latent Terms: Sparse BM25-Style Retrieval over Sparse Autoencoder Vocabularies}.
  arXiv:2605.29384, 2026.

\bibitem{unstablefeatures2026}
  (Anonyme).
  \textit{Unstable Features, Reproducible Subspaces}.
  arXiv:2606.12138, 2026.

\bibitem{identifiablesae2026}
  (Anonyme).
  \textit{Toward Identifiable Sparse Autoencoders}.
  arXiv:2605.31245, 2026.

\bibitem{multilingualconcepts2025}
  (Anonyme).
  \textit{Sparse Autoencoders Can Capture Language-Specific Concepts Across Diverse Languages}.
  arXiv:2507.11230, 2025.

\bibitem{beckmannqueloz2026}
  Beckmann, P. \& Queloz, M.
  \textit{Mechanistic Indicators of Understanding in Large Language Models}.
  2026.

\end{thebibliography}

\end{document}
